\documentclass[10pt,a4paper]{amsart}
\usepackage[margin=19mm]{geometry}
\usepackage{amsmath,amssymb,mathtools}
\usepackage[scr=rsfs,cal=euler]{mathalfa}
\usepackage{xfrac}
\usepackage[unicode,hidelinks,bookmarks=false]{hyperref}
\newcommand{\ie}{i.e.\ }
\newcommand{\cf}{cf.\ }
\newcommand{\resp}{respectively\ }
\newcommand{\Subsection}{\subsection}
\providecommand{\nobreakdash}{\nobreak\hbox{-}}
\setlength{\emergencystretch}{2em}
\multlinegap0pt
\usepackage{bm}
\usepackage{bbm}
\usepackage{dsfont}
\usepackage{xspace}
\usepackage{cancel}
\usepackage{mathtools}
\usepackage{xcolor}
\usepackage{amsthm}
\makeatletter
\@ifundefined{theorem}{\newtheorem{theorem}{Theorem}}{}
\@ifundefined{definition}{\newtheorem{definition}{Definition}}{}
\@ifundefined{lemma}{\newtheorem{lemma}[theorem]{Lemma}}{}
\makeatother
\title[Weak Correlations as the Underlying Principle for Linearizat]{Weak Correlations as the Underlying Principle for Linearization of Gradient-Based Learning Systems}
\author{Ori Shem-Ur, Khen Cohen, Aviv Orly, Yaron Oz}
\date{Source: arXiv:2401.04013v3 (CC BY 4.0)}
\begin{document}
\maketitle

\noindent\emph{Source and licence.} Weak Correlations as the Underlying Principle for Linearization of Gradient-Based Learning Systems. Version arXiv:2401.04013v3 (CC BY 4.0), distributed under the Creative Commons Attribution 4.0 International licence, as recorded in the arXiv licence metadata for this exact version and shown on the abstract page https://arxiv.org/abs/2401.04013v3. Full rights notice: licenses/arxiv-2401.04013-CC-BY-4.0.txt.\par\smallskip

\noindent\textit{Editorial scope.} Counted unit: the Lemma whose source label is \texttt{lem:TenMagEquivalency} in the article (source0.tex lines 1315--1333, source label \texttt{lem:TenMagEquivalency}), delivered together with the article's own complete original proof of it (source0.tex lines 1356--1424, one \texttt{proof} environment of 69 lines). No mathematics is written, repaired or re-derived: statement and proof are the article's own text line for line. The proof is detached from the statement in the source: the article states the result at lines 1315--1333 and prints its proof at lines 1356--1424, and the proof environment's own header names the statement, so the association is the article's own. Required context retained uncounted: def:BoundTenMag (lines 670--684). Every definition, display and label of those lines is retained unchanged. Omitted from this package and not redistributed: 1--669, 685--1314, 1334--1355, 1425--3003. Nothing inside a retained excerpt is omitted or altered: every retained body is delivered exactly as the article's lines are written, so no omission receipt of any kind accompanies this package. Citations: the retained excerpts carry no citation command at all, so no bibliography entry of the article is transcribed and no reference is redistributed. Reference adaptation: none was needed and none was made; every reference inside the delivered unit resolves inside it, so no unresolved reference or citation is produced. Printed slips kept literal: no printed word, symbol, hypothesis or number of the counted statement or of its proof was altered. Layout and class: the article's own class is \texttt{article} with packages including iclr2026\_conference, times, amsmath, amsfonts, bm, hyperref, url, inputenc, fontenc, booktabs, nicefrac, microtype, xcolor, graphicx; the wrapper is the lane's pinned \texttt{amsart} editorial wrapper at 10pt on A4, which restates the theorem-like environments the retained text needs and adds the article's own macro definitions that the retained text needs (none), with extra wrapper packages (bm, bbm, dsfont, xspace, cancel, mathtools, xcolor). The delivered numbering is the unit's own. Assets: no retained span contains a figure environment, a graphics-inclusion command, a PDF-page inclusion, a table or a TikZ picture; no image file is copied into this package, the package declares no asset dependency and no image file is redistributed with it. Licence: version arXiv:2401.04013v3 of this article is distributed under the Creative Commons Attribution 4.0 International licence, as recorded in the arXiv licence metadata for this exact version and shown on the abstract page https://arxiv.org/abs/2401.04013v3. A full rights notice is delivered at licenses/arxiv-2401.04013-CC-BY-4.0.txt. Characters: every retained excerpt is pure ASCII, so the delivered engine, XeLaTeX with its default Latin Modern text font, reports no missing character.
\begin{definition}
[Asymptotic Upper Bound of Random Tensors]
\label{def:BoundTenMag}
\textup{A random tensor \(M\), as defined above, is said to be asymptotically upper bounded by \(f\in\mathcal{N}\) if:
\begin{equation}
M= O\left(f\right) \ ,
\end{equation}
if and only if:
\begin{equation}
\forall g\in\mathcal{N}\,\text{s.t.}\,f=o\left(g\right):\lim_{n\rightarrow\infty}P\left(\left\Vert M_n\right\Vert \leq g\left(n\right)\right)=1\ .
\label{eq:BoundTenMag1}
\end{equation}
The lower asymptotic bound, \(f=\Omega (M)\), is defined analogously, but with the inequality reversed and \(g=o(f)\).
}
\end{definition}

\begin{lemma}[Equivalent Definitions for Tensor's Asymptotic Bound]
\label{lem:TenMagEquivalency}
\textup
{For any random tensor \(M\) and \(f\in\mathcal{N}\), the two definitions for bounding the tensor's asymptotic behavior \(O(M)\leq O(f)\) are equivalent (the first is the original definition, (\ref{def:BoundTenMag})):
\begin{enumerate}
\item 
\begin{equation}
\forall g\in\mathcal{N}\,\text{s.t.}\,f=o\left(g\right):\lim_{n\rightarrow\infty}P\left(\left\Vert M_{n}\right\Vert \leq g\left(n\right)\right)=1 \ .
\label{eq:MagDefBound1}
\end{equation}
\item
\begin{equation}
\lim_{c\rightarrow\infty}\lim_{n\rightarrow\infty}P\left(\left\Vert M_{n}\right\Vert \leq cf\left(n\right)\right)=1 \ .
\label{eq:MagDefBound2}
\end{equation} 
\end{enumerate}
(The same applies for \(O(f)\leq O(M)\)).
    }
\end{lemma}

\begin{proof}[\textbf{Proof - Lemma \ref{lem:TenMagEquivalency}}]
\

We will prove the two directions of the lemma separately.

\underline{Assuming the second condition in \textcolor{black}{Equation} \ref{eq:MagDefBound2} is satisfied:}

Given some \(0<p<1\), we know using \textcolor{black}{Equation} \ref{eq:MagDefBound2} that there is some \(0<c\) such that for sufficiently large \(n\in\mathbb{N}\):
\begin{equation}
p\leq P\left(\left\Vert M_n\right\Vert \leq cf\left(n\right)\right) \ .
\end{equation}
Given some \(g\in\mathcal{N}\) such that \(f=o(g)\), we know that for sufficiently large \(n\in\mathbb{N}\):
\begin{equation}
cf\left(n\right)\leq g\left(n\right) \ ,
\end{equation}
which means that for sufficiently large \(n\in\mathbb{N}\):
\begin{equation}
p\leq P\left(\left\Vert M_n\right\Vert \leq cf\left(n\right)\right)\leq P\left(\left\Vert M_n\right\Vert \leq g\left(n\right)\right) \ .
\end{equation}
As we proved that for any \(0<p<1\) we get that:
\begin{equation}
\lim_{n\rightarrow\infty}\left(P\left(\left\Vert M_n\right\Vert \leq g\left(n\right)\right)\right)=1 \ .
\end{equation}
And as we proved that for any arbitrary \(g\in\mathcal{N}\) such that \(f=o(g)\), we proved the first part of the lemma.

\underline{Assuming the first condition in \textcolor{black}{Equation} \ref{eq:MagDefBound1} is satisfied:}

If we assume in contradiction that \textcolor{black}{Equation} \ref{eq:MagDefBound2} is not satisfied, we get that there is some \(0<p<1\) such \textcolor{black}{that}:
\begin{equation}
\forall n_{0}\in\mathbb{N},\,0<c,\,\exists n_{0}\leq n\in\mathbb{N}:P\left(\left\Vert M_n\right\Vert \leq cf\left(n\right)\right)<p \ .
\end{equation}
In particular, \textcolor{black}{this} means that if we choose the sequence \(\left\{ c_{i}=i\right\} _{i=1}^{\infty}\), there are \(\tilde{n}_1<\tilde{n}_2<\tilde{n}_3\ldots\in\mathbb{N}\) such \textcolor{black}{that}:
\begin{equation}
    \forall i\in\mathbb{N}:P\left(\left\Vert M_{\tilde{n}_{i}}\right\Vert \leq if\left(\tilde{n}_{i}\right)\right)<p \ .
    \label{eq:Zin:Flin1}
\end{equation}
The reason that we can require that \(\left\{ \tilde{n}_{i}\right\} _{i=1}^{\infty}\) is \textcolor{black}{increasing} is that we know that we can find such \(n\)-\textcolor{black}{values} for any sufficiently large \(n_0\) and for any \(c\). So by induction we can require every time that every \(\tilde{n}_i\) is bigger than all previous \(\tilde{n}\)-\textcolor{black}{values}.

\underline{\textcolor{black}{Assuming the second condition of Equation \ref{eq:MagDefBound2} is satisfied:}}

Suppose, by contradiction, that \textcolor{black}{Equation} \ref{eq:MagDefBound1} is not satisfied. Then, there exists some \(0 < p < 1\) such that:
\begin{equation}
\forall n_0 \in \mathbb{N},\, 0 < c,\, \exists n_0 \leq n \in \mathbb{N}: P\left(\left\Vert M_n\right\Vert \leq cf\left(n\right)\right) < p \ .
\end{equation}
In particular, if we choose the sequence \(\forall i \in \mathbb{N}: c_i = i\), there exist \(\tilde{n}_1 < \tilde{n}_2 < \tilde{n}_3 \dots \in \mathbb{N}\) such that:
\begin{equation}
\forall i \in \mathbb{N}: P\left(\left\Vert M_{\tilde{n}_i}\right\Vert \leq if\left(\tilde{n}_i\right)\right) < p \ .
\end{equation}
Since we can find such \(n\)-values for any sufficiently large \(n_0\) and any \(c\), \textcolor{black}{we} can require by induction that each \(\tilde{n}_i\) is greater than all previous \(\tilde{n}\)-values. 

We can now define the function:
\begin{equation}
\forall n \in \mathbb{N}: g\left(n\right) = \left(\max\left\{ i \in \mathbb{N} \mid \tilde{n}_i \leq n\right\}\right) f\left(n\right) \ .
\label{eq:Zin:Flin2}
\end{equation}
Since \(\left\{ \tilde{n}_i\right\}_{i=1}^{\infty}\) is increasing, we know by the Archimedean property that \(\max\left\{ i \in \mathbb{N} \mid \tilde{n}_i \leq n\right\}\) is also increasing and unbounded, which implies:
\begin{equation}
\lim_{n \rightarrow \infty} \frac{g\left(n\right)}{f\left(n\right)} = \lim_{n \rightarrow \infty} \max\left\{ i \in \mathbb{N} \mid \tilde{n}_i \leq n\right\} = \infty \ .
\end{equation}
However, by using \textcolor{black}{Equations} \ref{eq:Zin:Flin1} \textcolor{black}{and} \ref{eq:Zin:Flin2}, we also have:
\begin{equation}
    \forall n_0 \in \mathbb{N},\, \exists n_0 \leq n \in \mathbb{N}: P\left(\left\Vert M_n\right\Vert \leq g\left(n\right)\right) < p \ ,
\end{equation}
which means that:
\begin{equation}
\lim_{n \rightarrow \infty} P\left(\left\Vert M_n\right\Vert \leq g\left(n\right)\right) \neq 1 \ .
\end{equation}
This contradicts our assumption in \textcolor{black}{Equation} \ref{eq:MagDefBound1}. Therefore, by reductio ad impossibile, \textcolor{black}{Equation} \ref{eq:MagDefBound2} must be satisfied, completing the proof for the second direction.
\end{proof}

\end{document}
